\begin{recipe}{Bloemkoolpasta}
\kicker[Hoofdgerecht,Vlees,Doordeweeks]{Hoofdgerecht · vlees · doordeweeks}
\index[register]{Bloemkoolpasta}

\meta{Voor 4 personen \dvd Snel \dvd Op gevoel}

\lettrine{R}{omige} pasta met geroosterde spekjes, zoete erwtjes en
geblancheerde bloemkool. Snel, troostrijk en flexibel. De hoeveelheden gaan
helemaal op gevoel.

\blockrule{Ingrediënten}
\begin{ingredients}
  \ing{1}{bloemkool, in roosjes}\index[register]{Bloemkool}
  \ing{1}{ui}\index[register]{Ui}
  \ing{1}{knoflookteentje}\index[register]{Knoflook}
  \ing{250 g}{spekblokjes}\index[register]{Spekjes}
  \ingb{scheut kookroom}
  \ingb{handje diepvrieserwtjes}\index[register]{Erwtjes}
  \ingb{pasta naar keuze}
  \ingb{parmezaan of pecorino}\index[register]{Parmezaan}
  \ingb{zout en versgemalen peper}
\end{ingredients}

\blockrule{Bereiding}
\tip{Bewaar altijd een kopje pastawater. Daarmee maak je de saus precies romig
genoeg.}
\begin{steps}
  \item Snipper de ui en knoflook en snijd de bloemkool in kleine roosjes: zo
        garen ze sneller en mengen ze beter door de pasta.
  \item Ontdooi de diepvrieserwtjes in een kommetje kokend water; giet af zodra
        ze op kamertemperatuur zijn.
  \item Blancheer de bloemkool kort: een paar minuten in kokend water, dan
        meteen afgieten. Zo wordt hij net zacht maar blijft hij nog een
        beetje stevig, in plaats van slap te koken.
  \item Kook de pasta 1 minuut korter dan de verpakking aangeeft; in de saus
        wordt hij nog wat zachter.
  \item Verhit een hapjespan en bak de spekjes. Voeg vlak voor ze bruin zijn de ui toe
        en bak mee.
  \item Is de ui glazig? Voeg de bloemkool toe en roerbak 1–2 minuten mee.
  \item Voeg de erwtjes en kookroom toe, zet het vuur laag en doe de deksel op
        de pan.
  \item Als alles warm is: deksel eraf, pasta erbij (bewaar een kopje pastawater).
        Rasp er Parmezaan of pecorino door, versgemalen peper naar smaak, en maak los met wat
        pastawater als de saus te dik is.
\end{steps}
\end{recipe}
